\section{Correlation lengths, pulling, and prescribed spatial endpoints}
\label{R:resp:section}

There are two further ways to change the thermal law. One may reward displacement by an exponential factor, or prescribe the terminal vertex. The first change tests convergence of a generating function. The second asks for one coefficient in two spatial coordinates. We first determine the exponential threshold and the associated force response. We then prove the spatial local estimate needed to normalize a prescribed endpoint.

In this section lengths $\ell$ count edges of ordinary walks from a fixed vertex $o$. Put $z_t=\rho e^{-t}$, $t>0$, and define
\[
 G_t(b)=\sum_{\gamma:o\to b}z_t^{\ell(\gamma)},\qquad
 \chi(t)=\sum_bG_t(b).
\]
The empty walk is included. For a unit vector $e$, let
\begin{align}
 m_e(t)&=\sup\left\{s\ge0:\sum_b e^{s e\cdot(b-o)}G_t(b)<\infty\right\},
 \label{R:resp:directional-mass}\\
 m_{\mathrm{rad}}(t)&=\sup\left\{s\ge0:\sum_b e^{s|b-o|}G_t(b)<\infty\right\}.
 \label{R:resp:radial-mass}
\end{align}
Both convergence sets are intervals containing zero. For the directional sum this follows by separating nonnegative and negative projections; the latter contribution is bounded by $\chi(t)$. All constants below allow the fixed conversion between Euclidean units and row spacing.

Directional decay and endpoint tilts have long been studied through renewal methods. Ioffe proved all-direction Ornstein--Zernike asymptotics for self-avoiding walk on $\mathbb Z^d$ in the strictly subcritical fugacity regime~\cite{Ioffe1998}. Borgs, Chayes, King and Madras studied directional masses for anisotropic walks and linear displacement of the most likely endpoint under asymmetric positive step weights~\cite{BorgsChayesKingMadras2000}. Ioffe and Velenik proved ballistic local limit results for self-interacting walks under fixed forces in the ballistic regime~\cite{IoffeVelenik2008}. Our estimates concern the approach to zero force and critical fugacity on the honeycomb lattice. We derive the required uniform powers from the quantitative bridge estimates above.

\begin{theorem}[Exponential thresholds and force]\label{R:resp:force}
Uniformly over unit vectors $e$, as $t\downarrow0$,
\[
 m_e(t)=t^{3/4+o(1)},\qquad m_{\mathrm{rad}}(t)=t^{3/4+o(1)}.
\]
More precisely, for every $\varepsilon>0$ and all sufficiently small $t$,
\[
 t^{3/4+\varepsilon}\le m_{\mathrm{rad}}(t)\le m_e(t)\le C t^{3/4},
\]
where $C$ is independent of $e$. For each unit vector $e$ and $s\ge0$, the limit
\[
 f_e(s)=\lim_{j\to\infty}\frac1j\log
 \sum_{\ell(\gamma)=j}\rho^j e^{s e\cdot(\gamma_j-o)}
\]
exists. It satisfies $f_e(0)=0$ and
\[
 f_e(s)=s^{4/3+o(1)}\qquad(s\downarrow0),
\]
uniformly in $e$.
\end{theorem}

The proof uses two consequences of the absolute count and fast-path bounds in Propositions~\ref{R:geo:absolute-count} and~\ref{R:geo:fast}. For every fixed sufficiently small $\xi>0$,
\begin{equation}
 \sum_{\substack{R\le D(\gamma)\le2R\\ L(\gamma)\le R^{4/3-\xi}}}
       \rho^{L(\gamma)}\le C_\xi e^{-R^{c_\xi}},
 \qquad
 U_R:=\sum_{D(\gamma)\le R}\rho^{L(\gamma)}
       \le e^{C_\nu(1+R)^\nu}\quad(\nu>0).
 \label{R:resp:critical-inputs}
\end{equation}
These sums start at a specified vertex and count vertices in $L$. For the first bound, paths of diameter at least $R$ have length at least $cR$. Choose $\eta>0$ so small that $(4/3-\xi)(3/4+\eta)<1$. At each remaining length $n\le R^{4/3-\xi}$, the refined fast-path estimate applies and gives mass at most
\[
 \exp(-c_\eta R/n^{3/4})\le\exp(-c_\eta R^{3\xi/4}).
\]
Summing the polynomial number of lengths gives the first bound. For the second, a self-avoiding trace of diameter at most $R$ visits at most $C(1+R)^2$ lattice vertices. Apply the absolute count bound with exponent smaller than $\nu/2$ and sum over these lengths. Changing either sum to edge weights costs one constant factor.

For the divergence estimate we need a positive mass of bridges whose thermal discount is bounded below. Theorem~\ref{R:renew:finite-input} gives $B_h\ge ch^{-1/4}$ and $M_h\le Ch^{13/12}$. Hence
\[
 B_h\{L>C_0h^{4/3}\}\le \frac{M_h}{C_0h^{4/3}}
       \le \frac C{C_0}h^{-1/4}.
\]
Choosing the fixed $C_0$ large enough retains a positive fraction of $B_h$ with $L\le C_0h^{4/3}$. This uses the finite first-length mass before any exact-length conditioning.

\begin{proof}[Proof of Theorem~\ref{R:resp:force}]
Fix $\varepsilon>0$ and put $R=t^{-3/4-\varepsilon}$. Stop a walk on its first visit to distance at least $R$ from its starting vertex. Its endpoint lies at distance at most $R+a_{\mathrm{lat}}$, where $a_{\mathrm{lat}}$ is the lattice edge length, and its diameter lies between $R$ and $2R+a_{\mathrm{lat}}$. Let $W(t,R)$ be the total $z_t$-weight of these stopped prefixes. The prefixes with at most $R^{4/3-\xi}$ vertices have mass tending to zero by the first estimate in~\eqref{R:resp:critical-inputs}, with harmless fixed adjustments of the radii. The remaining prefixes have total mass at most
\[
 C U_{3R}\exp(-tR^{4/3-\xi}).
\]
Choose $\xi>0$ so small that
\[
 4/3-\xi-\frac1{3/4+\varepsilon}>0,
\]
and then choose $\nu$ smaller than this difference in the bound for $U_{3R}$. It follows that $W(t,R)=o(1)$.

Now split any complete walk greedily into such stopped prefixes, restarting the rule at the endpoint of each prefix, followed by one remainder whose displacement is less than $R$. Edge weights multiply. Dropping avoidance between different pieces gives, at $s=1/R$,
\[
 \sum_\gamma z_t^{\ell(\gamma)}e^{sA(\gamma)}
 \le e\chi(t)\sum_{k\ge0}
       \bigl(e^{1+a_{\mathrm{lat}}/R}\max_{\text{vertex types}}W(t,R)\bigr)^k<\infty
\]
for all sufficiently small $t$. The estimates hold at both vertex types by lattice symmetry. Thus $m_{\mathrm{rad}}(t)\ge t^{3/4+\varepsilon}$, and $m_e(t)\ge m_{\mathrm{rad}}(t)$.

To prove divergence, take an integer height $H\asymp t^{-3/4}$. Choose one of the three inward port normals at $o$ whose scalar product with $e$ is at least $1/2$. Consider bridges in that direction with height $h\in[H,2H]$, length at most $C_1H^{4/3}$, and horizontal displacement of the sign favorable to $e$. Reflection in the normal preserves length and critical weight, so the total mass at every such height is at least $cH^{-1/4}$. Every retained displacement has projection on $e$ at least $c_0H$, where $c_0>0$ is independent of $e$.

From each retained bridge extract the prefix ending at its first renewal cut whose relative height lies in $[H,2H]$ and whose prefix has the stated sign and length properties. Such a cut exists because the terminal cut qualifies. A fixed extracted prefix may have suffixes to several heights in $[H,2H]$; their total critical mass is at most
\[
 \sum_{0\le j\le H}B_j\le CH^{3/4}.
\]
Since the original bridges summed over heights have mass at least $cH^{3/4}$, the distinct extracted prefixes have total critical mass at least a fixed $c_1>0$. Their lengths are at most $C_1H^{4/3}$, so their $e^{-tL}$-discounted mass is also bounded below by a fixed positive constant.

Concatenate any number of these prefixes in the chosen direction. The pieces lie in successive disjoint open slabs. Their joins are recoverable by the same first-qualifying-cut rule, relative to the current initial port. Later pieces cannot change an earlier cut or its prefix tests. Hence different ordered lists give different port paths. Trimming the initial and final half-edges gives ordinary plane walks from $o$; the length and endpoint-weight corrections are bounded independently of the number of pieces. Consequently the directionally tilted plane sum diverges when
\[
 c_2e^{s c_0H}>1.
\]
This proves $m_e(t)\le C/H\le Ct^{3/4}$, uniformly in $e$.

It remains to pass from convergence thresholds to the fixed-length free energy. Write
\[
 a_j(s,e)=\sum_{\ell(\gamma)=j}\rho^j
             e^{s e\cdot(\gamma_j-o)}.
\]
After an even number of steps the terminal vertex has the same lattice type as $o$. Splitting there and dropping avoidance therefore gives
$a_{2j+2k}\le a_{2j}a_{2k}$. Fekete's lemma gives a limit for $(2j)^{-1}\log a_{2j}$. Single-edge truncation gives
\[
 a_{j+1}\le 3\rho e^{s a_{\mathrm{lat}}}a_j.
\]
Applying this twice places each odd-length logarithm between the adjacent even-length logarithms up to a constant. Thus the limit defining $f_e(s)$ exists at all lengths. It is finite. The connective-constant identity gives $f_e(0)=0$.

The directional susceptibility is $\sum_j e^{-tj}a_j(s,e)$. The root test shows that it converges if $t>f_e(s)$ and diverges if $t<f_e(s)$. In particular, $s<m_e(t)$ implies $f_e(s)\le t$, and $s>m_e(t)$ implies $f_e(s)\ge t$. Applying the threshold bounds with $t=s^{1/(3/4+\delta)}$ and $t=s^{1/(3/4-\delta)}$, with an arbitrarily small fixed $\delta>0$ and harmless constant adjustments, proves $f_e(s)=s^{4/3+o(1)}$. All threshold constants used here are uniform in $e$.
\end{proof}

\begin{corollary}[One-sided force response]\label{R:resp:force-derivatives}
For $s>0$ the convex function $f_e$ has one-sided derivatives, and
\[
 f'_{e,-}(s)=s^{1/3+o(1)},\qquad f'_{e,+}(s)=s^{1/3+o(1)}
 \qquad(s\downarrow0).
\]
At a point where $f_e$ is differentiable, the mean displacement in direction $e$, divided by length, under the force-weighted $j$-edge law tends to $f'_e(s)$ as $j\to\infty$.
\end{corollary}
\begin{proof}
Each $j^{-1}\log a_j(s,e)$ is convex, and its pointwise limit is convex. Since $f_e(0)=0$, convexity gives
\[
 \frac{f_e(s)}s\le f'_{e,-}(s)\le f'_{e,+}(s)
       \le\frac{f_e(2s)-f_e(s)}s\le\frac{f_e(2s)}s.
\]
Here $f_e(s)\ge0$: rotation by $120$ degrees about $o$ makes the mean endpoint of the unforced fixed-length law zero, so Jensen's inequality gives $a_j(s,e)\ge a_j(0,e)$. Theorem~\ref{R:resp:force} now proves the derivative exponents. Finally the derivative of $j^{-1}\log a_j$ is the stated normalized mean. Its lower and upper difference quotients converge to those of $f_e$; letting their increments tend to zero proves the claim at differentiability points. No differentiability assertion is needed for the one-sided result.
\end{proof}

The exponent $1/3$ agrees with Pincus's force-extension scaling relation $1/\nu-1$ at $\nu=3/4$~\cite[Eq.~(I.4)]{Pincus1976}; here it follows from the threshold bounds and convexity. This small-force relation is distinct from fixed-force ballistic asymptotics; see also~\cite[Section~4.2.3]{IoffeVelenik2010}.

For prescribed endpoints we will use the same geometric estimates before summing over lengths. In edge-count notation, the absolute critical count bound in Proposition~\ref{R:geo:absolute-count} is
$c_j\rho^j\le\exp(C_\nu j^\nu)$ for every $\nu>0$, where $c_j$ counts rooted $j$-edge walks. The refined fast-path bound in Proposition~\ref{R:geo:fast} is
\begin{equation}
 \sum_{\substack{\ell(\gamma)=j\\D(\gamma)\ge d}}\rho^j
       \le\exp\bigl(-c_\eta d/j^{3/4}\bigr),
 \qquad d\ge j^{3/4+\eta},
 \label{R:resp:refined-fast}
\end{equation}
for fixed $\eta>0$ and sufficiently large $j$. The conversion from vertex count to edge count changes length by one and weight by a fixed factor; decreasing the slack and the exponential constant absorbs both changes. Both bounds concern unnormalized masses, so they remain available before a prescribed endpoint has been normalized.

\begin{remark}[A second route to the radial threshold]\label{R:resp:radial-alternative}
The turning-chain estimates give a useful alternative which requires no first-length estimate at a single chosen height. Write $p$ for the critical irreducible law, $T$ for its height and $L$ for its vertex length. By Lemma~\ref{R:geo:increment-tails} and the length deficit in Proposition~\ref{R:renew:increments},
\[
 p(e^{-L/N};T>R)\ge cR^{-3/4}-q_N,
 \qquad q_N:=p(1-e^{-L/N})=N^{-9/16+o(1)}.
\]
At $R=N^{3/4-\delta}$ the right side is at least $c'R^{-3/4}$. A pull $s=N^{-3/4+2\delta}$ in a fixed positive multiple of the height therefore makes the total tilted irreducible weight exceed one. Free concatenation then forces divergence of the radially tilted plane sum, after the bounded endpoint conversion.

For convergence put $s=N^{-3/4-\delta}$ and choose an even block length $j\asymp N^{1+\eta}$ with $\eta$ sufficiently small after $\delta$. On $D\le j^{3/4+\eta}$ the pull is bounded, while $e^{-j/N}$ beats the critical count bound. On successive larger diameter ranges, \eqref{R:resp:refined-fast} beats the pull because $s\ll j^{-3/4}$. The tilted mass of this block is less than one at both vertex types. Splitting paths into such blocks, with one bounded-length remainder, proves convergence. This recovers $m_{\mathrm{rad}}(1/N)=N^{-3/4+o(1)}$ by an argument directly adapted to the absolute turning-chain bounds.
\end{remark}

\subsection{A spatial local limit with lattice compatibility}
The exponential threshold alone gives no lower bound for $G_t(b)$ at a specified vertex. We now prove that a sum of order $H^{3/4}$ irreducible bridges can reach every admissible target of height $H$ at the expected two-dimensional lattice scale. The proof does not assume convergence to a unique stable law.

Let $Y=(T,X)$ be the displacement of one $p$-distributed irreducible, with $T$ measured in row spacings and $X$ in units of the spacing between ports on one cut. The ports on row $h$ are shifted horizontally by $h/2$ modulo this spacing. Thus their relative-displacement lattice is
\[
 \Gamma=\{(h,m+h/2):h,m\in\mathbb Z\}.
\]
The two elementary two-vertex bridges have displacements $(1,1/2)$ and $(1,-1/2)$. They are irreducible and generate $\Gamma$, so $\Gamma$ is also the lattice generated by all possible increments. Let $\Delta$ be the Euclidean diameter of an irreducible and $L$ its vertex length. Put $\beta=3/4$ and $K=H^\beta$.

For Fourier inversion the relevant subgroup is $\Gamma_1$, generated by differences of possible increments. Lemma~\ref{R:geo:atoms} proves that $\Gamma_1$ has rank two and hence finite index in $\Gamma$. All increments have one common class $a_*+\Gamma_1$; that class generates $\Gamma/\Gamma_1$. We call $k$ compatible with $y\in\Gamma$ when $ka_*-y\in\Gamma_1$.

\begin{proposition}[Prescribed displacement with controlled length and diameter]
\label{R:resp:point-local}
For each $C<\infty$ and $\eta>0$ there are $c>0$, $M_0<\infty$ and $H_0<\infty$ such that, for every integer $H\ge H_0$ and relative-port target $y=(H,x)\in\Gamma$ with $|x|\le CH$,
\begin{equation}
 \sum_{\lceil K\rceil\le k\le\lfloor2K\rfloor}
 p^{\otimes k}\left\{
 \sum_{i=1}^kY_i=y,\quad
 \sum_{i=1}^kL_i\le H^{4/3+\eta},\quad
 \sum_{i=1}^k\Delta_i\le M_0H
 \right\}
 \ge c K H^{-2}=cH^{-5/4}.
 \label{R:resp:point-bound}
\end{equation}
Incompatible counts contribute zero. In particular the paths counted here remain within distance $M_0H$ of their initial port.
\end{proposition}

\begin{proof}
We first establish the estimate without the length and diameter restrictions. Although subsequential laws of the rescaled sums may differ, each has full forward support and a strictly positive density in the open forward half-plane. Fourier inversion will then give a uniform lattice lower bound. Finally we mark individual increments to impose the two restrictions.

\emph{Subsequential jump limits.} Lemma~\ref{R:geo:increment-tails} gives
\begin{equation}
 p(|Y|>r)\le Cr^{-\beta},\qquad
 p(|Y|;|Y|\le r)\le Cr^{1-\beta}.
 \label{R:resp:jump-bounds}
\end{equation}
Consider any sequence $H\to\infty$. On a subsequence the measures
\[
 \mu_H=H^\beta p(Y/H\in\cdot)
\]
converge vaguely on $\mathbb R^2\setminus\{0\}$ to a measure $\mu$. This follows by compactness on each closed annulus and a diagonal extraction. The bounds
\[
 \mu_H(|y|>r)\le Cr^{-\beta},\qquad
 \int_{|y|\le r}|y|\,d\mu_H(y)\le Cr^{1-\beta}
\]
give the corresponding bounds for $\mu$, by approximation away from zero. In particular $\int(1\wedge|y|)\,d\mu(y)<\infty$. The measure is supported in $\{T\ge0\}$.

If $k/H^\beta\to u>0$, expansion of one characteristic function gives
\begin{equation}
 \mathbb E\exp\left(is\cdot H^{-1}\sum_{i=1}^kY_i\right)
 \longrightarrow
 F_u(s):=\exp\left(u\int(e^{is\cdot y}-1)\,d\mu(y)\right).
 \label{R:resp:poisson-transform}
\end{equation}
For completeness, truncate the integral to $\delta\le|y|\le M$. Vague convergence gives the limit there at continuity radii. The small-jump error is at most $C|s|\delta^{1-\beta}$; the large-jump error is at most $CM^{-\beta}$. Send first $H$ to infinity, then $\delta$ to zero and $M$ to infinity. The one-factor error is $O(H^{-\beta})$, so its squared term disappears after multiplication by $k$.

The right side is a characteristic function. Construct independent Poisson point collections on disjoint annuli, with intensity $u\mu$, and sum their points. There are finitely many points outside the unit disk, and the sum of the norms inside it has finite expectation. Thus the small-jump series converges absolutely almost surely. This is the finite-variation form of the Poisson construction of a pure-jump infinitely divisible law; see Schilling~\cite[Chapter~7]{Schilling2016} for the general construction. Tightness of the rescaled sums also follows directly from~\eqref{R:resp:jump-bounds}: remove jumps above $M$, and use the truncated first moment for the others. Fourier uniqueness therefore proves weak convergence to the constructed law. We prove the support and uniform lattice lower bound needed here next.

\emph{Full forward support.} For any $u>0$ the support of the limiting law is the closed additive semigroup generated by the support of $\mu$, including zero, and is independent of $u$. One inclusion follows by truncating the Poisson sum. Conversely, a specified finite list of jumps can be approximated with positive probability, excluding every other jump above a small cutoff. The remaining small jumps have arbitrarily small total norm with positive probability, by their first-moment bound. Hence each finite sum of support points belongs to the support.

This semigroup contains the entire open forward half-plane. To prove it, fix an open ball $B$ whose closure lies in $\{T>0\}$ and then a smaller closed ball $B'$ in $B$. The free-terminal guide estimate in Section~\ref{R:geo:finite-inputs}, applied over a fixed positive interval of target heights, gives
\begin{equation}
 \sum_{k\ge0}p^{*k}(HB')\ge c_B H^\beta.
 \label{R:resp:bridge-ball}
\end{equation}
Indeed a narrow straight guide places the terminal port in $HB'$, uniformly over order $H$ heights, and contributes $cH^{-1/4}$ at each height. Counts $k\le\xi H^\beta$ have total contribution at most $C_B\xi^2H^\beta$: to reach $HB'$ their height sum must exceed $c_BH$, and the truncated-sum bound from~\eqref{R:resp:jump-bounds} gives probability at most $C_B kH^{-\beta}$. Counts $k\ge S H^\beta$ contribute at most $C_BH^\beta e^{-c_BS}$, because
\[
 p\left(\sum_{i=1}^kT_i\le C_BH\right)
 \le e^{C_B}\bigl(p(e^{-T/H})\bigr)^k
 \le C_Be^{-c_B k/H^\beta}.
\]
Choose $\xi$ small and $S$ large. The remaining mass in~\eqref{R:resp:bridge-ball} is at least $c'_BH^\beta$, spread over at most $SH^\beta+1$ counts. Thus for some $k_H\in[\xi H^\beta,SH^\beta]$, the probability of $H^{-1}\sum_{i\le k_H}Y_i\in B'$ is bounded below independently of $H$. On a further subsequence $k_H/H^\beta\to u\in[\xi,S]$. Weak convergence and the closedness of $B'$ show that the limiting law gives positive mass to $B'$. Because all positive $u$ have the same support, every such ball meets the support of every limiting law.

\emph{Fourier inversion on the displacement lattice.} We have obtained full support in the open forward half-plane. To turn this into a positive lattice lower bound, we need a density and Fourier domination. Write $\varphi(\theta)=p(e^{i\theta\cdot Y})$ for the increment characteristic function. Lemma~\ref{R:geo:atoms} gives, near each of its modulus-one characters, the uniform estimate
\begin{equation}
 |\varphi(\theta)|\le 1-c\operatorname{dist}(\theta,\mathcal P)^\beta,
 \label{R:resp:fourier-decay}
\end{equation}
where $\mathcal P$ is the finite character group annihilating $\Gamma_1$ on the dual torus of $\Gamma$. Away from fixed neighborhoods of $\mathcal P$ the modulus is at most a constant strictly smaller than one. Passing to the limit near zero in~\eqref{R:resp:fourier-decay} gives $|F_u(s)|\le e^{-cu|s|^\beta}$. Thus the limiting law has a continuous density $g_u$.

In fact $g_u(v)>0$ whenever the first coordinate of $v$ is positive. The convolution identity $g_u=g_{u/2}*g_{u/2}$ follows from~\eqref{R:resp:poisson-transform}. Choose a small ball of heights strictly between zero and the first coordinate of $v$. Full support and continuity provide a point in that ball where $g_{u/2}>0$, and a smaller ball on which it is bounded below. The reflected difference ball about $v$ remains in the forward half-plane and has positive mass under the other half-time law. The convolution is therefore positive at $v$.

Take compatible $k$ and $y$ with $k/H^\beta\to u\in[1,2]$ and $y/H\to v$. Fourier inversion on $\Gamma$ yields
\begin{equation}
 H^2p^{*k}(y)\longrightarrow c_\Gamma g_u(v),
 \qquad c_\Gamma>0.
 \label{R:resp:lattice-limit}
\end{equation}
Here $c_\Gamma$ is the covolume of $\Gamma_1$ in the chosen coordinates. To check the periodicity factor, for $\theta_0\in\mathcal P$ all increments have the same phase $e^{i\theta_0\cdot a_*}$, so
$\varphi(\theta_0+\theta)=e^{i\theta_0\cdot a_*}\varphi(\theta)$. Compatibility makes the phase in the inversion integrand equal to one. Thus all the finitely many neighborhoods contribute identically. In each neighborhood rescale $\theta$ by $H^{-1}$ and dominate by $e^{-c'|s|^\beta}$ using~\eqref{R:resp:fourier-decay}; the complement is exponentially small in $k$. This proves~\eqref{R:resp:lattice-limit} with its stated constant.

The same subsequence argument proves a uniform positive lower bound when $|x|/H\le C$ and $k/H^\beta\in[1,2]$ is compatible. If it failed, choose offending targets and counts, extract $x/H\to v_2$, $k/H^\beta\to u$, and then a vague limit $\mu$ as above. Equation~\eqref{R:resp:lattice-limit} would contradict $g_u(1,v_2)>0$. A fixed positive fraction of the integers in $[K,2K]$ are compatible with each $y\in\Gamma$, because $\Gamma/\Gamma_1$ is finite cyclic. Consequently
\begin{equation}
 \sum_{K\le k\le2K}p^{*k}(y)\ge c_CKH^{-2}.
 \label{R:resp:unrestricted-point}
\end{equation}

\emph{Length and diameter restrictions.} We now remove paths whose length or total increment diameter is too large. If $\sum_iL_i>M$, then
\[
 \sum_i(1-e^{-L_i/M})\ge1-e^{-\sum_iL_i/M}\ge1-e^{-1}.
\]
Mark one increment by $1-e^{-L/M}$. Its total marked mass is $q_M:=p(1-e^{-L/M})$. The other $k-1$ increments match the remaining displacement with probability at most $C(k-1)^{-8/3}\le CH^{-2}$ by Lemma~\ref{R:geo:atoms}, independently of the value of the marked increment. Summing the mark and then the count gives the excluded mass bound
\[
 C K^2H^{-2}q_M.
\]
For $M=H^{4/3+\eta}$, the weak deficit bound $q_M=M^{-9/16+o(1)}$ makes this $o(KH^{-2})$.

The same argument controls spatial excursions without a new independence assumption. If $\sum_i\Delta_i>BH$, then $\sum_i\min(1,\Delta_i/(BH))\ge1$. The diameter tail implies
\[
 p\bigl(\min(1,\Delta/(BH))\bigr)\le C(BH)^{-\beta}.
\]
Marking by this function excludes mass at most $CB^{-\beta}KH^{-2}$. Choose a fixed $B=M_0$ large enough to leave at least half the lower bound in~\eqref{R:resp:unrestricted-point}, and then take $H$ large enough for the length exclusion. The entire concatenated path is within its sum of increment diameters of its first port. This proves~\eqref{R:resp:point-bound}.
\end{proof}

\subsection{Thermal walks ending at a specified vertex}
We now use the local estimate to construct actual paths in a corridor between two prescribed vertices. This supplies a lower denominator before any conditional tail estimate is used.

\begin{theorem}[Prescribed spatial offsets]\label{R:resp:prescribed}
At each fixed $N>0$ put $t=1/N$. Define
\[
 \underline m(N)=\liminf_{|b-o|\to\infty}
        \frac{-\log G_{1/N}(b)}{|b-o|},\qquad
 \overline m(N)=\limsup_{|b-o|\to\infty}
        \frac{-\log G_{1/N}(b)}{|b-o|},
\]
where both limits range over all lattice vertices and directions. Then
\[
 \underline m(N)=N^{-3/4+o(1)},\qquad
 \overline m(N)=N^{-3/4+o(1)}\qquad(N\to\infty).
\]
If a sequence of terminal vertices satisfies $|b_N-o|=N^{3/4+o(1)}$, the probability law
\[
 \mathbb P_{N,b_N}(\gamma)=
       \frac{(\rho e^{-1/N})^{\ell(\gamma)}}{G_{1/N}(b_N)},
       \qquad \gamma:o\to b_N,
\]
satisfies $\ell=N^{1+o(1)}$ and $D=N^{3/4+o(1)}$ in probability. For each fixed $r>0$ it also satisfies
\[
 \mathbb E_{N,b_N}\ell^r=N^{r+o(1)},\qquad
 \mathbb E_{N,b_N}D^r=N^{3r/4+o(1)}.
\]
\end{theorem}

\begin{proof}
Fix a small $\delta>0$, and put $H=\lfloor N^{3/4-\delta}\rfloor$. Choose one of the three inward port normals at $o$ whose projection of $b-o$ is at least $c|b-o|$. Call the vertex type of $o$ type $P$, and the other type $R$. Start the port path on the cut immediately below $o$. If $b$ has type $R$, end at its incident port on the cut immediately above $b$; delete the initial and terminal half-edges to obtain the ordinary walk. If $b$ has type $P$, end at its incident port immediately below $b$; delete the initial half-edge and append the final half-edge from that port to $b$. The appended vertex lies beyond the final open slab and is therefore unvisited. If the port path has $L$ internal vertices, the resulting edge count is $L-1$ in the first case and $L$ in the second. These operations are recoverable from the fixed orientation and terminal vertex. The associated port displacement differs from $b-o$ by a bounded amount, has positive height comparable to $|b-o|$, and has horizontal component at most a constant times that height.

When $|b-o|\gg H$, divide this displacement along its straight guide into $q=O(|b-o|/H)$ exact port displacements, each with integer height between $H$ and $2H$ and bounded horizontal-to-height ratio. Round intermediate positions to ports on the chosen cuts; this changes each horizontal displacement by a bounded amount and leaves the final port exact. Proposition~\ref{R:resp:point-local}, with a fixed length slack $\eta>0$ sufficiently small after $\delta$, gives critical mass at least $cH^{-5/4}$ for every segment, using only lengths at most $(2H)^{4/3+\eta}$ and diameter at most $C H$. Since
\[
 H^{4/3+\eta}=o(N),
\]
the thermal discount on each segment is bounded below by a fixed positive constant.

Concatenate the segments in their consecutive slabs. They cannot share vertices across slabs. The predetermined intermediate cuts are each crossed once, so the pieces are recoverable from their concatenation. They remain in a corridor of width $CH$ about the straight guide. After the bounded endpoint conversion, we obtain
\begin{equation}
 G_{1/N}(b)\ge
       \exp\left[-C_\delta\left(1+\frac{|b-o|}{H}\right)\log N\right]
 \qquad(|b-o|\gg H).
 \label{R:resp:denominator}
\end{equation}

For fixed $N$ and $s<m_{\mathrm{rad}}(1/N)$,
\[
 G_{1/N}(b)\le
 \left(\sum_v e^{s|v-o|}G_{1/N}(v)\right)e^{-s|b-o|}.
\]
The parenthesized quantity is finite and independent of $b$. Hence the lower spatial rate is at least $m_{\mathrm{rad}}(1/N)$, which is at least $N^{-3/4-\varepsilon}$ for every fixed $\varepsilon>0$ and sufficiently large $N$. On the other hand~\eqref{R:resp:denominator} gives the upper spatial rate at most $C_\delta H^{-1}\log N$. Letting the fixed exponent slacks tend to zero proves both asserted rate scales. No equality or existence of a single directional spatial limit is required.

Now suppose $|b_N-o|=N^{3/4+o(1)}$. For every fixed $\zeta>0$, choose $\delta<\zeta/3$ in~\eqref{R:resp:denominator}. For all sufficiently large $N$ this gives
\begin{equation}
 G_{1/N}(b_N)\ge e^{-N^\zeta}.
 \label{R:resp:subexponential-denominator}
\end{equation}
The choice of $\zeta$ will be made after the desired tail slack, and stays fixed as $N$ tends to infinity.

Let $c_j\rho^j$ be the total critical mass at edge length $j$. Proposition~\ref{R:geo:absolute-count} gives $c_j\rho^j\le\exp(C_\nu j^\nu)$ for every $\nu>0$. Thus, for every $\varepsilon>0$, the absolute thermal mass on $j>N^{1+\varepsilon}$, even multiplied by $j^r$ for any fixed $r$, is at most $\exp(-cN^\varepsilon)$ after decreasing $\nu$. Dividing by~\eqref{R:resp:subexponential-denominator} proves the length upper tail and its moment version.

For the spatial upper tail choose a small $\varepsilon'>0$ relative to $\varepsilon$. On lengths $j\le N^{1+\varepsilon'}$, the event $D>N^{3/4+\varepsilon}$ is covered by~\eqref{R:resp:refined-fast}, with a sufficiently small fixed $\eta$ there. Its absolute mass is at most $\exp(-N^c)$ for some $c>0$, after summing the polynomial number of lengths. The lengths above this cutoff have just been treated. Choose $\zeta<c$ in~\eqref{R:resp:subexponential-denominator}. This proves the spatial upper tail. The same proof, using $D\le a_{\mathrm{lat}}j$ and the weighted length tail, gives its positive-moment version.

Finally every path to $b_N$ has $D\ge|b_N-o|$. If $j<N^{1-\varepsilon}$, this required displacement is at least $N^{3/4-o(1)}$, so~\eqref{R:resp:refined-fast} again bounds the absolute mass by $\exp(-N^{c'})$ for some $c'>0$; lengths smaller than the graph distance contribute nothing. Division by~\eqref{R:resp:subexponential-denominator} proves the length lower tail. The diameter lower bound is deterministic. Integrating the proved weighted upper tails and using these lower bounds gives both moment exponents.
\end{proof}

\begin{corollary}[Length responses at a prescribed endpoint]\label{R:resp:thermal-derivatives}
For each fixed vertex $b$ and $t>0$, $G_t(b)$ can be differentiated term by term any fixed number of times. If $|b_N-o|=N^{3/4+o(1)}$, then for each positive integer $j$,
\[
 \left.\frac{(-1)^j\partial_t^jG_t(b_N)}{G_t(b_N)}\right|_{t=1/N}
       =N^{j+o(1)},\qquad
 -\left.\partial_t\log G_t(b_N)\right|_{t=1/N}=N^{1+o(1)}.
\]
The terminal vertex is held fixed while taking each derivative.
\end{corollary}
\begin{proof}
On a compact subinterval of $t>0$, every differentiated series is dominated by a constant times the untilted susceptibility at a smaller positive value of $t$. Hence
$(-1)^j\partial_t^jG_t(b)/G_t(b)$ is exactly the $j$th edge-length moment under the normalized endpoint law. Apply Theorem~\ref{R:resp:prescribed}. The logarithmic derivative is the first moment. Higher logarithmic derivatives are cumulants and are not identified by this argument.
\end{proof}
